\section{Mười hai nghĩa vụ tuân thủ và cách sinh lịch}

\boncau
{Bộ \soRule{} bản ghi khai báo trong mã (\texttt{COMPLIANCE\_RULES}): mã, tiêu đề,
tần suất, hạn chót, \texttt{applies\_to}, \texttt{legal\_refs} trỏ Điều trong corpus.}
{DNNVV cần nhắc mốc thuế, BHXH, lao động, BCTC. LLM không được suy hạn nộp ---
dễ bịa ngày.}
{Bộ sinh đọc hồ sơ tổ chức, tạo task cửa sổ 3 tháng trước + 12 tháng tới.
\texttt{INSERT \ldots ON CONFLICT DO NOTHING} theo (công ty, rule, kỳ).}
{Lúc đăng ký/cập nhật hồ sơ và khi user bấm sinh lại. Không chạy khi chưa gắn
tổ chức. Luật đổi hạn: sửa rule trong mã, không chờ mô hình nhớ.}

Mỗi mục dưới đây: nghĩa vụ là gì, công thức hạn trong mã, và khi nào rule được
lọc vào hồ sơ.

\subsection{VAT\_MONTHLY --- khai nộp GTGT theo tháng}

Áp dụng khi hồ sơ khai \texttt{vat\_period = monthly}. Hạn: ngày 20 của tháng
liền sau tháng phát sinh, theo hướng Điều 44 Luật Quản lý thuế (hồ sơ khai) và
Điều 55 (thời hạn nộp tiền) số 38/2019/QH14. \texttt{due\_day} $= 20$,
\texttt{due\_month\_offset} $= 1$.

Doanh nghiệp nhỏ thường được chọn quý; rule tháng tồn tại vì một số đơn vị vẫn
kê khai tháng, và vì demo cần thấy khác biệt khi đổi hồ sơ.

\subsection{VAT\_QUARTERLY --- khai nộp GTGT theo quý}

Áp dụng \texttt{vat\_period = quarterly}. Hạn được mã hoá ngày cuối tháng đầu
quý sau (\texttt{due\_day} $= 31$ kết hợp offset một tháng, bộ sinh kẹp theo số
ngày thật của tháng). Căn cứ: Điều 44 Luật Quản lý thuế và Điều 8 Nghị định
126/2020/NĐ-CP. Đây là rule người dùng DNNVV gặp nhiều nhất trên lịch.

\subsection{CIT\_PROVISIONAL --- tạm nộp thuế TNDN theo quý}

Mọi tổ chức trong phạm vi đồ án. Hạn ngày 30 tháng đầu quý sau. Rule nhấn:
không phát sinh hồ sơ khai tạm tính quý, chỉ nghĩa vụ tiền. Căn cứ Điều 55 Luật
Quản lý thuế và Điều 8 Nghị định 126/2020/NĐ-CP. Nhầm lẫn thường gặp là coi tạm
nộp như một tờ khai --- mô tả rule ghi rõ để UI không gọi là ``khai quý TNDN''.

\subsection{CIT\_FINALIZATION --- quyết toán TNDN năm}

Hạn ngày cuối tháng thứ ba sau khi kết thúc năm dương lịch
(\texttt{due\_fixed\_month} $= 3$, \texttt{due\_day} $= 31$). Căn cứ Điều 44 Luật
Quản lý thuế và Điều 12 Luật Thuế TNDN 14/2008/QH12 (kho còn số hiệu này; luật
mới cùng chủ đề được nạp riêng, rule chưa tự chuyển). Năm tài chính khác dương
lịch chưa tham số hoá --- hạn chế phạm vi.

\subsection{PIT\_FINALIZATION --- quyết toán TNCN năm}

Chỉ khi \texttt{min\_employees} $\ge 1$. Tổ chức trả thu nhập quyết toán thay
người lao động, hạn cùng cửa sổ tháng thứ ba. Căn cứ Điều 44 Luật Quản lý thuế
và Điều 24 Luật Thuế TNCN 04/2007/QH12. Công ty không có lao động không thấy
task này.

\subsection{BUSINESS\_LICENSE\_FEE --- lệ phí môn bài}

Hạn 30 tháng 01 hằng năm. Căn cứ Điều 5 Nghị định 139/2016/NĐ-CP và Điều 10
Nghị định 126/2020/NĐ-CP. Rule không tính số tiền theo vốn điều lệ --- đồ án
nhắc hạn, không thay phần mềm thuế.

\subsection{SOCIAL\_INSURANCE\_MONTHLY --- đóng BHXH hằng tháng}

Khi có ít nhất một lao động. Hạn ngày cuối tháng đang diễn ra
(\texttt{due\_month\_offset} $= 0$, \texttt{due\_day} $= 31$ kẹp cuối tháng).
Căn cứ Điều 34 Luật BHXH 41/2024/QH15 và Điều 4 Nghị định 58/2020/NĐ-CP (quỹ
tai nạn lao động). Tỷ lệ đóng không nằm trong rule; người dùng hỏi tỷ lệ thì
đi pipeline RAG, không đi lịch.

\subsection{LABOUR\_REPORT\_H1 và H2 --- báo cáo tình hình lao động}

Hai rule cùng căn cứ Điều 12 Bộ luật Lao động 45/2019/QH14 và Điều 4 Nghị định
145/2020/NĐ-CP, khác tháng: trước 05/06 và trước 05/12. Lọc
\texttt{min\_employees} $\ge 1$. Tách H1/H2 thay vì một rule ``hai lần/năm''
để task hiện thành hai mốc trên lịch tháng, dễ đánh dấu độc lập.

\subsection{LABOUR\_SAFETY\_REPORT --- báo cáo an toàn vệ sinh lao động}

Khi \texttt{min\_employees} $\ge 10$. Hạn 10 tháng 01 năm sau. Căn cứ Điều 81
Luật An toàn, vệ sinh lao động 84/2015/QH13. Ngưỡng 10 lao động là điều kiện
lọc trong mã, phản ánh hướng tập trung DNNVV có quy mô đủ để nghĩa vụ báo cáo
này trở nên thường gặp; khi văn bản đổi ngưỡng phải sửa \texttt{applies\_to}.

\subsection{FINANCIAL\_STATEMENT --- báo cáo tài chính năm}

Mọi tổ chức. Hạn mã hoá cuối tháng 3 (90 ngày sau 31/12 với năm dương lịch).
Căn cứ Điều 29 Luật Kế toán 88/2015/QH13 và Điều 80 Thông tư 133/2016/TT-BTC
(chế độ kế toán DNNVV). Nộp cơ quan thuế và thống kê được nêu trong mô tả;
hệ thống không tạo hai task tách cơ quan.

\subsection{INVOICE\_USAGE\_CHECK --- rà soát hóa đơn điện tử}

Rule quý, hạn ngày 15 tháng đầu quý sau: đây là \textbf{mốc nội bộ} nhắc rà
soát việc lập, gửi, lưu hóa đơn theo Nghị định 123/2020/NĐ-CP Điều 9 và Thông
tư 78/2021/TT-BTC Điều 6, không phải một tờ khai riêng có số hiệu thống nhất
như GTGT. Mục đích: kéo người dùng từ lịch sang tra cứu/hỏi đáp khi sắp đến
kỳ, chứ không tự đối soát dữ liệu tổng cục thuế.

\subsection{Bảng tổng hợp rule}

\begin{table}[htp]
\centering
\caption{\soRule{} rule tuân thủ: tần suất và điều kiện lọc.}
\label{tab:12-rule}
\begin{tabular}{p{3.6cm}p{2.2cm}p{2.4cm}p{5.2cm}}
\toprule
\textbf{Mã} & \textbf{Tần suất} & \textbf{Lọc hồ sơ} & \textbf{Hướng hạn trong mã} \\
\midrule
VAT\_MONTHLY & Tháng & Kỳ GTGT tháng & Ngày 20 tháng sau \\
VAT\_QUARTERLY & Quý & Kỳ GTGT quý & Cuối tháng đầu quý sau \\
CIT\_PROVISIONAL & Quý & Không & Ngày 30 tháng đầu quý sau \\
CIT\_FINALIZATION & Năm & Không & Cuối tháng 3 \\
PIT\_FINALIZATION & Năm & Có NLĐ & Cuối tháng 3 \\
BUSINESS\_LICENSE\_FEE & Năm & Không & 30 tháng 01 \\
SOCIAL\_INSURANCE\_MONTHLY & Tháng & Có NLĐ & Cuối tháng \\
LABOUR\_REPORT\_H1 & Năm & Có NLĐ & 05 tháng 6 \\
LABOUR\_REPORT\_H2 & Năm & Có NLĐ & 05 tháng 12 \\
LABOUR\_SAFETY\_REPORT & Năm & $\ge$ 10 NLĐ & 10 tháng 01 năm sau \\
FINANCIAL\_STATEMENT & Năm & Không & Cuối tháng 3 \\
INVOICE\_USAGE\_CHECK & Quý & Không & Ngày 15 tháng đầu quý sau \\
\bottomrule
\end{tabular}
\end{table}

Người dùng đổi số lao động hoặc kỳ GTGT rồi bấm sinh lại: task tương lai theo
rule mới xuất hiện; task quá khứ đã đánh dấu xong giữ nguyên. Không dùng LLM để
``đoán'' hạn nộp --- sai một ngày là sai nghiệp vụ.
